(1) 自旋二分之一粒子在 \(t=0\) 处于 \(\sigma_z\) 的上本征态 \(|\uparrow_z\rangle\)。随后哈密顿量突然变为 \(H=(\hbar\Omega/2)(\sin\theta\,\sigma_x+\cos\theta\,\sigma_z)\)，其中 \(\Omega>0\)、\(0\le\theta\le\pi\)。求精确演化算符 \(U(t)\)。
(2) 求在 \(t>0\) 测量 \(\sigma_z\) 得到向下的概率 \(P_{\downarrow}(t)\)。
(3) 求随时间的 \(\langle\sigma_z\rangle\) 和哈密顿量期望值，说明为什么磁场不做时变功。
(4) 检查 \(\theta=0,\pi/2,\pi\) 与 \(\Omega t=2\pi\) 时的极限，并区分全局自旋相位与物理测量概率。
